\section{接する直線・二次曲線から生じる十八次元成分}
\label{sec:tangent-component-v33}

相異なる直線と二次曲線が接する場合には、横断的に交わる場合とは別の
混合方向が現れる。そのうち一つは全 Poisson スキームの新たな成分を与える。
この節でも重みは $(1,1,1,3)$ である。

\begin{theorem}[接線型の成分]
\label{thm:tangent-component-v33}
次の点を取る。
\begin{equation}
 c_*=x(xz-y^2),\qquad
 D_*=-\frac52(x\partial_y+2y\partial_z),\qquad
 F_*=(0,\tfrac12yz^3,z^4).
 \label{eq:tangent-point-v33}
\end{equation}
$(c_*,D_*)$ の $\operatorname{SL}_3$ 軌道上で混合核の全体を動かした族の閉包は、
十八次元の既約成分 $\mathcal E_{\mathrm{tan}}$ である。
この成分は一般点で滑らかであり、重み付き座標変換群の稠密な一軌道を持つ。
従って既知の八成分とは相異なり、重み $(1,1,1,3)$ の Poisson スキームには
少なくとも九つの既約成分が存在する。
\end{theorem}

\begin{proof}
まず $D_*(c_*)=0$ であり、直接代入により $M(c_*,D_*;F_*)=0$ を得る。
$c_*$ を消す発散零の線形場は
$\mathbb C(x\partial_y+2y\partial_z)$ に限られる。
この場は $D_*$ とも可換なので、組 $(c_*,D_*)$ の
$\operatorname{SL}_3$ 安定化群の Lie 代数は一次元である。
従ってその軌道 $\mathcal O$ は滑らかな既約七次元多様体である。
この安定化群の計算は、三次式の十係数に関する線形方程式だけで確認できる。
下記の検証では軌道の接写像の七次小行列式 $12$ も与えている。

混合行列 $M(c_*,D_*;-)$ は Euler 商の三十五次元空間から $S_6$ への写像であり、
階数は二十四、核の次元は十一である。
この階数計算は後述の重みブロック表から従う。
同変性により $\mathcal O$ 全体で階数は一定なので、その核は階数十一のベクトル束を成す。
従ってその全空間は滑らかで既約、次元は $7+11=18$ である。

次に全 Jacobi 方程式の一次化を調べる。変数は
$(\dot c,\dot D,\dot F)\in S_3\oplus\mathfrak{sl}_3\oplus(S_4^{\oplus3}/xS_3)$ とし、
一次化は
\begin{align*}
 &D_*(\dot c)+\dot D(c_*),\\
 &2F_*(\dot c)-\dot c\operatorname{div}F_*
  -(x\times\dot D)\cdot F_*+M(c_*,D_*;\dot F)
\end{align*}
である。これは $38\times53$ 行列であり、厳密な有理数演算により階数三十五を持つ。
具体的な小行列式は
\begin{equation}
 -\frac{5958618044776701215625}{128}\ne0
 \label{eq:tangent-jac-minor-v33}
\end{equation}
である。よって点 \eqref{eq:tangent-point-v33} での Zariski 接空間は十八次元となる。
この点を通る上記十八次元族があるため、局所次元も十八であり、同点は全スキームの滑らかな点である。
その族の閉包はしたがって既約成分であり、一般点で滑らかである。

稠密軌道も明示できる。この点の二次場の係数は
\[
 p=\left(\frac43y^2-\frac16xz,-\frac73xy,\frac76x^2\right).
\]
三次式による $w$ の移動の混合項は $p\times\nabla g$、$g\in S_3$ である。
$p$ の三成分の最大公約因子は一であり、$\operatorname{curl}p=D_*\ne0$ なので、
$p\times\nabla g=0$ なら $\nabla g=ap$、$a\in\mathbb C$ である。
さらに curl を取ると $a=0$ となり、$g=0$ である。
Euler 商においても核は生じない。実際
$p\times\nabla g=h x$ なら $p$ との内積から $h c_*=0$、従って $h=0$ となる。
従って移動の像は十次元である。

対角場 $H=x\partial_x-z\partial_z$ の重みを使うと、$\Pi$ の重みは一であり、
$p\times\nabla g$ の最小重みは $1-3=-2$ である。
一方 $F_*$ の重みは $-3$ であり、第一成分が零で他が非零なので Euler 方向ではない。
従って Euler 商でも非零であり、移動の像に含まれない。
従って全混合核は
\[
 [F]=[p\times\nabla g]+u[F_*],\qquad g\in S_3,\quad u\in\mathbb C
\]
と一意に書ける。$u\ne0$ の場合、三次式による $w$ 移動で $g$ を消し、
$w$ の非零定数倍によって $u$ を一にできる。
基底の $\operatorname{SL}_3$ 軌道も合わせると、$u\ne0$ の部分全体が一つの座標変換軌道となる。
最後に、従来の八成分の次元は全て十九以上なので、この十八次元成分と一致しない。
\end{proof}

\begin{proposition}[二次元 Lie 代数による代表]
\label{prop:tangent-affine-lie-v33}
点 \eqref{eq:tangent-point-v33} の重み付きEuler切断による発散零の持ち上げは
\[
 \widehat P_*=H\wedge N,\qquad
 H=\tfrac14(3x\partial_x+y\partial_y-z\partial_z-5w\partial_w),\qquad
 N=x\partial_y+2y\partial_z-z^3\partial_w
\]
であり、
\[
 [H,N]=\tfrac12N,\qquad \operatorname{div}H=-\tfrac12,
 \qquad\operatorname{div}N=0
\]
を満たす。
\end{proposition}

\begin{proof}
一般の持ち上げの式に $\operatorname{div}F_*=9z^3/2$ と $D_*$ を代入すると、
三変数側の係数は
$(y^2/2+xz/4,-3xy/2,3x^2/4)$、混合係数は
\[
 (-3xz^3/4,\,-yz^3/4+5xw/4,\,z^4/4+5yw/2)
\]
となる。これらは $H\wedge N$ の係数そのものである。
括弧と発散の恒等式も各座標に作用させれば従う。
\end{proof}

この代表は、非可換な二次元 Lie 代数の多項式作用から得られる。
さらに $N$ は三角的な局所冪零導分である。
この幾何学的記述は、前の成分認定と独立に、量子化を構成するための具体的な入口を与える。

\begin{proposition}[明示的な有理第一積分]
\label{prop:tangent-first-integral-v33}
\[
 q=xz-y^2,\qquad
 J=x^4w+x^3yz^3-2x^2y^3z^2+\frac85xy^5z-\frac{16}{35}y^7
\]
と置く。$q,J$ の重みはそれぞれ $2,7$ であり、
\[
 N(q)=N(J)=0,\qquad H(q)=\tfrac12q,\qquad H(J)=\tfrac74J
\]
となる。従って $J^2/q^7$ は $X$ 上の有理関数として定義され、
$\mathcal E_{\mathrm{tan}}$ の代表の有理第一積分となる。

対応する水平一形式を
\[
 \alpha=c_*\,dw-w\,dc_*+w(x\times D_*)\cdot dx+(F_*\times x)\cdot dx
\]
とすると、正確な多項式恒等式
\begin{equation}
 2q\,dJ-7J\,dq=2x^3\alpha
 \label{eq:tangent-pencil-v33}
\end{equation}
を持つ。特に
\[
 d(J^2/q^7)=\frac{2x^3J}{q^8}\,\alpha.
\]
\end{proposition}

\begin{proof}
$N$ と $H$ を各単項式に作用させれば最初の四つの等式が得られる。
式 \eqref{eq:tangent-pencil-v33} は $dx,dy,dz,dw$ の四係数を比較すれば従う。
最後の式は商の微分から得られる。
\end{proof}

従ってこの成分の稠密軌道は有理的に可積分である。
式 \eqref{eq:tangent-pencil-v33} は重み $(2,7)$ の二つの多項式から作る
特殊な pencil の微分形式から共通因子 $x^3$ を除いた表示である。
次数九の形式をそのまま用いる通常の $R(2,7)$ 成分と、この次数六の成分とを
同一視するものではない。

\subsection{計算の再現可能な指定と二つの追加共鳴}

計算に使う単項式は
\[
 x^iy^jz^{d-i-j},\qquad i=0,\ldots,d,\quad j=0,\ldots,d-i
\]
の順とする。$\mathfrak{sl}_3$ の基底は、まず $x_j\partial_i$ を
$i=1,2,3$、$j=1,2,3$、$i\ne j$ の順に取り、最後に
$x\partial_x-y\partial_y$、$x\partial_x-z\partial_z$ を加える。
$F$ の基底は成分順、単項式順に取り、第一成分の $x$ で割れる単項式を除く。
行は $D(c)$ の三次係数、続いて $M(c,D;F)$ の六次係数とする。
この順序で零始まりの添字を使うと、式 \eqref{eq:tangent-jac-minor-v33} の行集合は
$\{0,\ldots,37\}\setminus\{0,10,11\}$、列集合の補集合は
\[
 \{12,15,18,19,23,28,38,39,40,41,43,44,45,46,47,48,50,51\}
 \subset\{0,\ldots,52\}
\]
である。これにより巨大な行列全体を印刷せずに小行列式を完全に指定できる。

より一般に $c=x(xz-y^2)$、$D=a(x\partial_y+2y\partial_z)$ と置く。
$H=x\partial_x-z\partial_z$ の重みによって混合行列を分解すると、
標的重み $k$ のブロックは次の通りである。
表の $g_k(a)$ は、そのブロックの最大階数の全小行列式の monic な最大公約式である。
空の二ブロック $k=-6,-5$ は省略した。
\[
\begin{array}{c|c|c|l}
 k&\text{行数}\times\text{列数}&\text{一般階数}&g_k(a)\\\hline
 -4&2\times1&1&1\\
 -3&2\times2&2&(a-2)(a+1)\\
 -2&3\times3&3&(a-2)(a+1)(a+\tfrac52)\\
 -1&3\times4&3&(a-2)(a+1)\\
 0&4\times5&4&(a-2)(a+1)(a+\tfrac32)\\
 1&3\times5&3&a-2\\
 2&3\times5&3&a-2\\
 3&2\times4&2&a-2\\
 4&2\times3&2&a-2\\
 5&1\times2&1&a-2\\
 6&1\times1&1&a-2
\end{array}
\]
全小行列式の最大公約式を求めているため、これは有限個の値で階数を試した表ではなく、
複素パラメータ $a$ 全体での階数降下点を尽くす証明書である。
各根での降下は一ブロックにつき一であり、混合核は
\[
 \dim\ker M=
 \begin{cases}
 20&a=2,\\
 14&a=-1,\\
 11&a=-\tfrac52,-\tfrac32,\\
 10&\text{それ以外}
 \end{cases}
\]
となる。$a=-\tfrac32$ については混合核の増大だけから成分性を結論しない。
その方向の既存成分への所属、または新たな成分の有無は残る問題である。
\texttt{verify\_tangent\_component\_v33.py} は、表の全小行列式、
軌道の接写像、二次場・混合核の恒等式、全Jacobi行列の小行列式を有理数上で検証する。
